\section{Đối tượng và phạm vi nghiên cứu}
Đối tượng nghiên cứu của khóa luận là hệ thống giám sát và điều khiển môi trường trong nhà sử dụng cảm biến, truyền thông không dây, nền tảng IoT và ứng dụng di động. Hệ thống được định hướng cho căn hộ chung cư, nơi số lượng node không quá lớn nhưng yêu cầu độ ổn định, dễ lắp đặt và dễ sử dụng. Cây trầu bà được xem là một thành phần hỗ trợ cải thiện vi khí hậu và giảm một phần chất ô nhiễm trong điều kiện phù hợp, còn hệ thống IoT giữ vai trò đo lường, điều khiển và cung cấp dữ liệu cho người dùng.

Phạm vi nghiên cứu tập trung vào một node cảm biến và một gateway. Node được đặt gần khu vực trồng trầu bà hoặc khu vực sinh hoạt cần giám sát. Các thông số đo gồm TDS, pH, nhiệt độ, độ ẩm, CO2 và ánh sáng. Phần điều khiển gồm quạt, đèn qua relay và điều hòa thông qua phát lệnh hồng ngoại. Khóa luận không đi sâu vào mô hình sinh học chi tiết của cơ chế hấp thụ VOC, không đánh giá định lượng tốc độ hấp thụ từng chất ô nhiễm và không thay thế các thiết bị đo IAQ chuyên dụng. Thay vào đó, đề tài tập trung vào xây dựng nguyên mẫu hệ thống có khả năng thu thập dữ liệu, cảnh báo và điều khiển thiết bị trong điều kiện căn hộ thực tế.
